\chapter{Introduction and Problem Definition}
\label{ch:intro}

\section{Project Background}
Hospital wards generate a continuous stream of vital signs and laboratory results, yet deterioration is
still often recognised late. Track-and-trigger scores such as the National Early Warning Score (NEWS)
\cite{smith2013news} and severity scores such as SAPS~II \cite{legall1993sapsii} reduce physiology to a
few thresholds, while machine-learning (ML) models trained on electronic health records (EHRs) report
higher discrimination \cite{rajkomar2018scalable,purushotham2018benchmarking}. At the same time, patients
rarely control who sees their own records, and clinicians who receive a patient from another facility
must often work without the history.

MediHaven is a web application, developed as a \CourseName\ in the \Department, that brings three
strands together in one teaching-scale system: (i) a physician triage dashboard that ranks synthetic
inpatients with four classical ML models and explains each vote, (ii) a patient-controlled medical vault
from which time-limited, scope-limited QR access passes can be shared, and (iii) an explicitly opt-in
research mode that rebuilds a real mortality-prediction task from the openly available MIMIC-IV Clinical
Database Demo \cite{johnson2023mimicivdemo,johnson2023mimiciv} and evaluates baselines with
leakage-aware, patient-grouped cross-validation.

MediHaven is an \textbf{educational software prototype}. It is not a medical device, it has not been
clinically validated, and none of its outputs may be used for diagnosis, treatment or other patient-care
decisions.

\section{Problem Statement}
Build a single, secure, testable web platform that (a) demonstrates explainable multi-model risk
stratification on realistic but synthetic inpatient data, (b) lets a patient grant a clinician
revocable, auditable and scope-restricted access to selected records, and (c) shows, on genuine
de-identified ICU data, how a mortality-risk pipeline must be built and evaluated honestly---without
leakage, with uncertainty, and without exposing patient-level data beyond the researcher's machine.

Existing approaches make this difficult: published models rarely release reproducible cohorts
\cite{johnson2017reproducibility}, deployed proprietary scores can perform far worse than reported
\cite{wong2021epicsepsis}, and record-sharing systems are typically institution-centred rather than
patient-controlled \cite{azaria2016medrec}.

\section{Motivation}
The project gives three groups something concrete: students and teachers obtain an end-to-end example
of clinical ML engineering (data pipeline, models, evaluation, serving, security and testing);
clinicians and evaluators can see how model votes, rule paths and uncertainty are presented; and
patients see a model of consent-based sharing with a visible access log. The research mode in
particular demonstrates the gap between an impressive demo accuracy on synthetic data and the much
more modest, uncertain performance obtained on real data (Chapter~\ref{ch:results}).

\section{Objectives}
\begin{enumerate}[label=O\arabic*.]
\item Implement role-based web portals for physicians, patients and administrators over a REST API.
\item Train and serve a four-model ensemble (K-Means, decision tree, KNN, MLP) that returns a risk tier,
      a score, per-model votes and an explanation for each synthetic patient.
\item Implement a patient vault with HMAC-signed, expiring, scope- and use-limited QR passes, an
      atomic use quota and an immutable access log.
\item Build a versioned MIMIC-IV Demo task (in-hospital mortality after the first 24 ICU hours) with
      explicit cohort, label and leakage rules, and evaluate baselines with patient-grouped
      cross-validation and bootstrap confidence intervals.
\item Serve verified research artifacts in a loopback-only research mode with integrity checks.
\item Protect secrets and data: strong-secret enforcement, CSP, rate limits, de-identified prompts
      to external AI, and no external AI in research mode.
\item Verify the system with an automated test suite and measure server performance under load.
\end{enumerate}

\section{Scope of the Project}
\textbf{Included:} the features above; a synthetic cohort of 100 patients; the MIMIC-IV Demo~2.2
(100 patients, 140 ICU stays); a production entry point (Waitress) and container recipe.
\textbf{Excluded:} real hospital integration (HL7/FHIR feeds), credentialed full MIMIC-IV, clinical
validation, real file uploads to the vault, and multi-instance deployment.
\textbf{Intended users:} students, teachers and evaluators; the physician/patient/administrator roles are
demonstration roles. \textbf{Operating environment:} Python~3.12 on Windows or Linux, a modern browser,
SQLite. \textbf{Expected scale:} a single server process serving tens of concurrent users
(Chapter~\ref{ch:scal}).

\section{Proposed Solution}
A layered Flask application (Chapter~\ref{ch:design}) serves Jinja2/JavaScript portals and a JSON API.
Synthetic predictions come from an in-memory ensemble whose weighted soft votes are fused into one of
four tiers; a decision-tree rule path and $k$ nearest historical cases make each result inspectable.
The vault issues HS256 JSON Web Tokens \cite{rfc7519} (HMAC \cite{rfc2104}) whose server-side record
holds the authoritative scope, expiry and use count. A separate research pipeline builds a hashed,
versioned dataset from MIMIC-IV Demo files outside the repository, trains five pipelines under grouped
cross-validation, and the web application serves only artifacts whose SHA-256 hashes verify.

\section{Novelty and Originality}
The individual algorithms are standard; the contribution is the \emph{integration and the engineering
discipline}: (i) one prototype combines explainable triage, patient-controlled sharing and a research
pipeline; (ii) research mode is fail-closed---it refuses unsafe secrets, non-loopback binding and
tampered artifacts, never falls back to synthetic data and never calls an external AI service;
(iii) the stay-level research view scores every ICU stay with an out-of-fold model and refuses to serve
unless the recomputed predictions reproduce the verified run exactly; and (iv) timestamp anomalies in the
outcome data are kept, flagged and analysed rather than silently dropped.

\section{Feasibility Analysis}
\textbf{Technical:} all components use mature open-source libraries (Flask, scikit-learn
\cite{pedregosa2011scikit}, pandas, DuckDB) and run on a laptop.
\textbf{Economic:} no licence costs; the MIMIC-IV Demo is open access; Gemini use is optional.
\textbf{Operational:} one command starts the demo (\texttt{run\_demo.py}); the production entry point
seeds and trains on first start. \textbf{Time and resources:} the repository history shows development
between 6 and 26 October 2026 (Section~\ref{sec:git}).

\section{Social Relevance and Practical Applications}
MediHaven is a teaching and demonstration artefact: it illustrates early-warning dashboards, consent-based
record sharing and responsible evaluation of clinical ML. Its practical value lies in education and in
prototyping user interfaces and data-protection patterns, not in patient care.

\section{Success Criteria and Key Performance Indicators}
Table~\ref{tab:kpi} lists the KPIs fixed for this project; Chapter~\ref{ch:testing} compares them with
the measured values.

\begin{table}[!htbp]
\centering
\caption{Key performance indicators and targets}
\label{tab:kpi}
\small
\begin{tabularx}{\textwidth}{@{}p{34mm}Yp{22mm}Y@{}}
\toprule
KPI & Definition & Target & Measurement method\\
\midrule
Functional correctness & passed / total automated tests & 100\,\% & \texttt{pytest} on the full suite\\
Response time & median latency of a risk prediction, 1 client & $\le$ 150 ms & timed HTTP requests to Waitress\\
Reliability & failed requests under 16 concurrent clients & 0 & 3{,}200-request load test\\
Throughput & simple API requests per second, 16 clients & $\ge$ 100 req/s & load test, \texttt{/api/vitals}\\
Resource use & server working set under load & $\le$ 512 MB & process-tree memory sampling\\
Synthetic model quality & hold-out accuracy on the synthetic split & $\ge$ 0.91 & project benchmark report\\
Research evaluation & discrimination with honest uncertainty & report AUROC with 95\% CI & grouped CV + bootstrap\\
\bottomrule
\end{tabularx}
\end{table}

\section{Organisation of the Report}
Chapter~\ref{ch:lit} reviews related systems and literature; Chapter~\ref{ch:srs} specifies requirements;
Chapter~\ref{ch:design} presents the architecture and design; Chapter~\ref{ch:algo} covers algorithms and
complexity; Chapter~\ref{ch:subjects} maps the work to the curriculum; Chapter~\ref{ch:impl} describes the
implementation; Chapters~\ref{ch:scal}--\ref{ch:results} report performance, testing and results; and
Chapter~\ref{ch:concl} concludes.
